% !TeX root = ../../../../../main.tex
% chapters/sections/chapter5/subsections/two_plane/equilibrium.tex

\section{自国金融政策の均衡への影響}
\label{sec:chapter5_two_plane_equilibrium}
本節では\((y_t^H,p_t^H)\)平面および\((c_t^{H\to W},p_t^{H\to W})\)平面を用いた2平面分析により、
自国金融政策が均衡にどのような影響を与えるかを簡潔に分析する。
\((y_t^H,p_t^H)\)平面においてはYD曲線およびYS-\(\hat{p}\)曲線を用い、
\((c_t^{H\to W},p_t^{H\to W})\)平面においてはCD曲線およびCS-\(\hat{p}\)曲線を用いる。
\(\beta\)ショックが起こってからの一連の動きについての分析は
\ref{sec:chapter5_beta_shock_interpretations}
でおこない、また\(a\)ショックが起こってからの一連の動きについての分析は
\ref{sec:chapter5_a_shock_interpretations}
でおこなう。

\subsection{\((y_t^H,p_t^H)\)平面における均衡への影響}
\label{sec:chapter5_two_plane_equilibrium_ysyd}
本モデルの均衡は\textbf{右上がりのYS-\(\hat{p}\)}
\eqref{form:chapter5_two_plane_ys_derivation_p_H_modified_eq}
と\textbf{右下がりのYD曲線}
\eqref{form:chapter5_two_plane_yd_derivation_p_H_eq}
の交点として描かれる。
\ref{sec:chapter5_two_plane_yd_preparation}
節で述べたようにYD曲線に含まれる\(X_t\)は自国金融政策の影響を受けない。
そのため自国金融政策の分析上考慮すべき変数は
\(i_t^H,\operatorname{E}_t[\lambda_{t+1}^H],\operatorname{E}_t[\hat{\pi}_{t+1}^H]\)
の3つである。
以下ではこれらを仮想的に単独で変化させ\((y_t^H,p_t^H)\)平面上の均衡の動きを観察する。

\subsubsection{名目利子率\(i_t^H\)の低下}
名目利子率\(i_t^H\)が低下するとYD曲線
\eqref{form:chapter5_two_plane_yd_derivation_p_H_eq}
の右辺が大きくなるため\textbf{YD曲線は右上へ移動}する。
これにより生産\(y_t^H\)および生産者物価指数\(p_t^H\)は上昇する。

\subsubsection{所得の期待限界効用\(\operatorname{E}_t[\lambda_{t+1}^H]\)の低下}
\label{sec:chapter5_two_plane_equilibrium_ysyd_E_lambda_H_fall}
所得の限界効用の期待\(\operatorname{E}_t[\lambda_{t+1}^H]\)が低下するとYD曲線
\eqref{form:chapter5_two_plane_yd_derivation_p_H_eq}
の右辺が大きくなるため\textbf{YD曲線は右上へ移動}する。
これにより生産\(y_t^H\)および生産者物価指数\(p_t^H\)は上昇する。
これは名目利子率\(i_t^H\)を下げる余地がないゼロ金利制約下にあっても、
中央銀行の約束が家計の期待に働きかけることにより生産\(y_t^H\)を回復させることが可能であることを示している。

\subsubsection{生産者物価指数（PPI）を用いたグロス・インフレ率の対数乖離の期待\(\operatorname{E}_t[\hat{\pi}_{t+1}^H]\)の上昇}
\label{sec:chapter5_two_plane_equilibrium_ysyd_E_pi_H_rise}
生産者物価指数（PPI）を用いたグロス・インフレ率の対数乖離の期待
\(\operatorname{E}_t[\hat{\pi}_{t+1}^H]\)が上昇した場合、YS-\(\hat{p}\)曲線
\eqref{form:chapter5_two_plane_ys_derivation_p_H_modified_eq}
の右辺が大きくなるため\textbf{YS-\(\hat{p}\)は左上へ移動}する。
これにより生産\(y_t^H\)は下落し、生産者物価指数\(p_t^H\)は上昇する。

\subsection{\((c_t^{H\to W},p_t^{H\to W})\)平面における均衡への影響}
\label{sec:chapter5_two_plane_equilibrium_cscd}
本モデルの均衡は\textbf{右上がりのCS-\(\hat{p}\)曲線}
\eqref{form:chapter5_two_plane_cs_derivation_p_H_W_eq}
と\textbf{右下がりのCD曲線}
\eqref{form:chapter5_two_plane_cd_derivation_p_H_W_c_H_W_modified_eq}
の交点として描かれる。
\ref{sec:chapter5_two_plane_cs_preparation}および
\ref{sec:chapter5_two_plane_cs_derivation}
節で述べたようにCS曲線に含まれる\(Z_t,\Omega_t\)は自国金融政策の影響を受けない。
そのため自国金融政策の分析上考慮すべき変数は
\(i_t^H,\operatorname{E}_t[\lambda_{t+1}^H],\operatorname{E}_t[\hat{\pi}_{t+1}^H]\)
の3つである。以下ではこれらを仮想的に単独で変化させ
\((c_t^{H\to W},p_t^{H\to W})\)平面上の均衡の動きを観察する。

\subsubsection{名目利子率\(i_t^H\)の低下}
\label{sec:chapter5_two_plane_equilibrium_cscd_i_H_fall}
名目利子率\(i_t^H\)が低下するとCD曲線
\eqref{form:chapter5_two_plane_cd_derivation_p_H_W_c_H_W_modified_eq}
の右辺が大きくなるため\textbf{CD曲線は右上へ移動}する。
これにより総消費指数\(c_t^{H\to W}\)および消費者物価指数\(p_t^{H\to W}\)は上昇する。

\subsubsection{所得の限界効用の期待\(\operatorname{E}_t[\lambda_{t+1}^H]\)の低下}
\label{sec:chapter5_two_plane_equilibrium_cscd_E_lambda_H_fall}
所得の限界効用の期待\(\operatorname{E}_t[\lambda_{t+1}^H]\)が低下するとCD曲線
\eqref{form:chapter5_two_plane_cd_derivation_p_H_W_c_H_W_modified_eq}
の右辺が大きくなるため\textbf{CD曲線は右上へ移動}する。
これにより総消費指数\(c_t^{H\to W}\)および消費者物価指数\(p_t^{H\to W}\)は上昇する。
これは名目利子率\(i_t^H\)を下げる余地がないゼロ金利制約下にあっても、
中央銀行の約束が家計の期待に働きかけることにより
総消費指数\(c_t^{H\to W}\)を回復させることが可能であることを示している。

\subsubsection{生産者物価指数（PPI）を用いたグロス・インフレ率の対数乖離の期待\(\operatorname{E}_t[\hat{\pi}_{t+1}^H]\)の上昇}
\label{sec:chapter5_two_plane_equilibrium_cscd_E_pi_H_fall}
生産者物価指数（PPI）を用いたグロス・インフレ率の対数乖離の期待
\(\operatorname{E}_t[\hat{\pi}_{t+1}^H]\)が上昇した場合、CS-\(\hat{p}\)曲線
\eqref{form:chapter5_two_plane_cs_derivation_p_H_W_eq}
の右辺が大きくなるため\textbf{CS曲線は左上へ移動}する。
これにより総消費指数\(c_t^{H\to W}\)は下落し、消費者物価指数\(p_t^{H\to W}\)は上昇する。